%! The Language of Variation — Unit 1, Lesson 1.1 — Homework (Answer Key)
% EFFL, AP Statistics variant: this is the lesson's SCORED practice and the LAST component in
% the packet. These items were the in-class Check Your Understanding set; they are now the
% graded homework, started in the closing minutes of the period and finished at home. The
% AP Practice component that precedes it stays assigned-but-unscored.
% \answerspace{H}{} reserves exactly H in the blank; the key passes the answer as the second
% argument, so the two files paginate identically by construction. Keep the macro and every
% \boxguard byte-identical in the blank and the key.%
% GENERATED FROM homework/main.tex — this file is the blank with answers filled in, so the
% \answerspace heights and every \boxguard are identical to it.
% NO teachernote here — teacher prose lives in the lesson plan.
\documentclass[10pt]{article}
\usepackage{apstats-article}
\usepackage{apstats-key}   % pulls in -boxes; do NOT also load -boxes
\newcommand{\answerspace}[2]{\par\nopagebreak\noindent\begin{minipage}[t][#1][t]{\linewidth}%
  \color{keyred}\bfseries #2\end{minipage}\par}

\begin{document}

\pageheader{Unit 1, Lesson 1.1}{Homework --- Answer Key}
% NAMESTRIP: no \namedateperiod here — the name row lives on the cover only.

\vspace{0.15in}

% Authored BYTE-IDENTICALLY in homework/ and homework_key/ — this box is what keeps the
% blank and the key the same number of pages.
\boxguard[8]
\begin{remindbox}
\small
\textbf{This is your graded homework.} It is scored and due next class. You will start it in the
last minutes of the period, so ask about anything that stalls you before you leave, then finish
it on your own.
\end{remindbox}

\vspace{0.12in}

% ── The scored practice (formerly this lesson's Check Your Understanding) ────
\boxguard[14]
\begin{notesbox}{Classifying Variables in a New Context}
Show the arithmetic wherever a question asks for a number, and write every explanation
\emph{in the context of the problem} --- that is what earns credit on the AP exam, and it is the
standard here.
\begin{enumerate}[leftmargin=1.9em, itemsep=7pt]
  \item A coffee shop logs every order it fills. Label each variable \textbf{C} or \textbf{Q}.

        \vspace{4pt}
        \renewcommand{\arraystretch}{1.3}
        \begin{tabularx}{\linewidth}{@{} Y Y Y Y @{}}
        \textbf{Drink size (S/M/L)} & \textbf{Number of shots} & \textbf{Price paid (\$)} &
        \textbf{Loyalty card \#} \\
        \ans{C} & \ans{Q} & \ans{Q} & \ans{C} \\
        \end{tabularx}

        The individuals here are \ans{the orders the shop filled}.
  \item \textbf{Same survey, two columns.} A pet survey records \emph{Number of pets owned} and
        \emph{Type of pet owned}. Classify each, then say what a summary of each column would
        look like.
        \answerspace{3.4cm}{Number of pets owned is quantitative (a count) --- summarize it with
        an average, such as 1.8 pets per household. Type of pet owned is categorical ---
        summarize it with counts or percents, such as 46\% dogs, 31\% cats.}
  \item A hospital records each patient's exact body temperature in $^\circ$C. It then adds a
        second column marking each patient \emph{Normal} or \emph{Fever}. Classify both
        columns, and explain where the second column came from.
        \answerspace{3.4cm}{Temperature is quantitative --- a measured amount. Normal/Fever is
        categorical --- a group label. The second column was built from the first by grouping the
        measurements at a cutoff, which trades amounts for labels and loses how far above the
        cutoff each patient was.}
  \item Name one variable that is written entirely in digits and is \textbf{not} quantitative,
        and say what its digits are doing instead of measuring.
        \answerspace{2.7cm}{A home ZIP code (also: jersey number, phone number, student ID). The
        digits identify which group a person belongs to rather than recording an amount, so
        averaging them is meaningless.}
  \item \textbf{AP-style multiple choice.} A survey of high school students records the five variables
        below. Circle the one that is \textbf{quantitative}, then justify it in one sentence.
        \begin{enumerate}[label=\textbf{(\Alph*)}, leftmargin=2.2em, itemsep=3pt]
          \item The number on the student's team jersey
          \item The number of the bus route the student rides
          \item \ans{The number of minutes the student spent on homework last night}
          \item The student's home ZIP code
          \item The student's cafeteria PIN
        \end{enumerate}
        \answerspace{2.7cm}{(C). Minutes are a measured amount, so 90 minutes really is twice
        45. The other four are labels written in digits.}
\end{enumerate}
\end{notesbox}

\vspace{0.12in}

% The next-lesson preview closes the packet, so it lives here rather than in AP Practice.
\boxguard[10]
\begin{spiralbox}
\small
\textbf{Next lesson: Representing a Categorical Variable.} You now know that a categorical
column is summarized by \emph{counting} rather than averaging. Next time we do that counting
properly --- frequency and relative frequency tables --- and see how the same counts can be
made to tell two very different stories.
\end{spiralbox}

\end{document}
